\subsection{导子的扩张}

命题 14. 设 A 为交换环，M 为 A-模，B 为 A-代数 \(\mathsf{T}(\mathbf{M})\) （相应地 \(\mathsf{S}(\mathbf{M})\) ，相应地 \(\bigwedge (\mathbf{M}))\) ），且 E 为一个 (B, B)-双模。设 \(d_0: \mathbf{A} \to \mathbf{E}\) 是环 A 到 A-模 E 中的一个导子，并设 \(d_{1}:M\to E\) 是一个加法群同态，使得对所有 \(a\in A\) 以及所有 \(x\in M\) ，

\[
d _ {1} (a x) = a d _ {1} (x) + d _ {0} (a). x,\tag{29}
\]

并且进一步，当 \(\mathbf{B} = \mathbf{S}(\mathbf{M})\) 时，

\[
x. d _ {1} (y) + d _ {1} (x). y = y. d _ {1} (x) + d _ {1} (y). x\tag{30}
\]

对所有 \(x, y\) （ \(\mathbf{M}\) 中的）成立，并且，当 \(\mathbf{B} = \bigwedge (\mathbf{M})\) 时，

\[
x. d _ {1} (x) + d _ {1} (x). x = 0\tag{31}
\]

对所有 \(x \in \mathbf{M}\) 成立。那么存在且仅存在一个导子 \(d\) （从 \(\mathbf{B}\) ——视为一个 \(\mathbf{Z}\) -代数——到 \((\mathbf{B}, \mathbf{B})\) -双模 \(\mathbf{E}\) ），使得 \(d|_{\mathbf{A}} = d_0\) 且 \(d|M = d_1\) 。

我们在 \(\mathbf{Z}\) -模 \(\mathbf{B} \oplus \mathbf{E}\) 上取由下式定义的结合 \(\mathbf{Z}\) -代数结构

\[
(b, t) (b ^ {\prime}, t ^ {\prime}) = (b b ^ {\prime}, b t ^ {\prime} + t b ^ {\prime})
\]

它以 (1, 0) 为单位元（第 8 号，命题 13）。在典范单射 \(t \mapsto (0, t)\) 下，E 等同于 B ⊕ E 的一个双边理想，使得 \(\mathrm{E}^2 = \{0\}\) 。另一方面，映射 \(h_0 \colon \mathbf{B} \oplus \mathbf{E}\) （由 \(h_0(a) = (a, d_0(a))\) 定义）是一个有单位元的环同态（第 8 号，命题 12）；在此映射下，B ⊕ E 于是成为一个 A-代数。此外，若对所有 \(x \in \mathbf{M}\) 记 \(h_1(x) = (x, d_1(x))\) ，则由 \(h_0\) 的定义以及 (29) 可知 \(h_1(ax) = h_0(a)h_1(x)\) ；换言之， \(h_1\) 是从 M 到 B ⊕ E 的一个 A-线性映射。于是存在且仅存在一个 A-代数同态 \(h \colon \mathbf{B} \to \mathbf{B} \oplus \mathbf{E}\) ，使得 \(h \mid \mathbf{M} = h_1\) （且必然有 \(h \mid \mathbf{A} = h_0\) ）：因为，若 \(\mathbf{B} = \mathsf{T}(\mathbf{M})\) ，这由 § 5 第 1 号的命题 1 得出；若 \(\mathbf{B} = \mathbb{S}(\mathbf{M})\) ，条件 (30) 表明 \(h(x)h(y) = h(y)h(x)\) 对 M 中所有 \(x, y\) 成立，结论由 § 6 第 1 号的命题 2 得出；最后若 \(\mathbf{B} = \bigwedge (\mathbf{M})\) ，条件 (31) 表明 \((h(x))^2 = 0\) 对所有 \(x \in \mathbf{M}\) 成立，又因 \(x \wedge x = 0\) ，结论由 § 7 第 1 号的命题 1 得出。同态 \(h\) 使得 \(p \circ h \colon \mathbf{B} \to \mathbf{B}\) 与增广 \(\rho \colon \mathbf{B} \oplus \mathbf{E} \to \mathbf{B}\) 的复合是恒等同态 \(l_B\) ，因为 \(p \circ h\) 与 \(l_B\) 在 A 的元素和 M 的元素上按定义一致，而这些元素组成的集合是 B 的一个生成系。因此我们可以写出 \(h(b) = (b, d(b))\) （对所有 \(b \in B\) ），且映射 \(b \mapsto d(b)\) （从 B 到 E 的）依第 8 号的命题 12 是一个具有所需性质的导子。

推论. 设 \(\mathbf{M}\) 是一个分次 \(\mathbf{K}\) -模（类型为 \(\Delta\) ）； \(\mathbf{K}\) -代数 \(\mathsf{T}(\mathbf{M})\) ， \(\mathsf{S}(\mathbf{M})\) 和 \(\bigwedge (\mathbf{M})\) 被赋予相应的类型为 \(\Delta' = \Delta \times \mathbf{Z}\) 的分次（ § 5，第 5 号，命题 7 ； § 6，第 6 号，命题 10 及 § 7，第 7 号，命题 11）。另一方面， \(\mathbf{M}\) 被赋予类型为 \(\Delta'\) 的分次，使得 \(\mathbf{M}_{\alpha,1} = \mathbf{M}_{\alpha}\) 对所有 \(\alpha \in \Delta\) 成立，且 \(\mathbf{M}_{\alpha,n} = \{0\}\) 对 \(\alpha \in \Delta\) 和 \(n \neq 1\) 成立。设 \(\varepsilon'\) 是 \(\Delta'\) 上的一个交换因子。

\begin{enumerate}
\def\labelenumi{(\roman{enumi})}
\item
  设 \(\mathbf{E}\) 是一个分次的（左和右） \(\mathsf{T}(\mathbf{M})\) -双模（类型为 \(\Delta'\) ）；对所有 \(\delta \in \Delta\) 以及每个整数 \(n \in \mathbf{Z}\) ，每个分次 \(\mathbf{K}\) -线性映射 \(f: \mathbf{M} \to \mathbf{E}\) （次数为 \(\delta_1' = (\delta, n)\) ）都可以唯一地扩张为一个 \(\varepsilon'\) -导子 \(d: \mathsf{T}(\mathbf{M}) \to \mathbf{E}\) （次数为 \(\delta'\) ）。
\item
  设 \(\mathbf{E}\) 是一个分次 \(\mathbb{S}(\mathbf{M})\) -模（类型为 \(\Delta'\) ）；为使分次 \(\mathbf{K}\) -线性映射 \(f: \mathbf{M} \to \mathbf{E}\) （次数为 \(\delta'\) ）可以扩张为一个 \(\varepsilon'\) -导子 \(d: \mathbb{S}(\mathbf{M}) \to \mathbf{E}\) （次数为 \(\delta'\) ），必要且充分条件是：对每个有序对 \((x, y)\) （由 \(\mathbf{M}\) 的齐次元素组成），
\end{enumerate}

\[
x. f (y) + \varepsilon_ {\delta^ {\prime}, (\deg (y), 1)} ^ {\prime} y. f (x) = y. f (x) + \varepsilon_ {\delta^ {\prime}, (\deg (x), 1)} ^ {\prime} x. f (y).\tag{33}
\]

该 \(\varepsilon'\) -导子 \(d\) 则是唯一的。

\begin{enumerate}
\def\labelenumi{(\roman{enumi})}
\setcounter{enumi}{2}
\tightlist
\item
  设 \(\mathbf{E}\) 是一个分次的（左和右） \(\bigwedge (\mathbf{M})\) -双模（类型为 \(\Delta'\) ）；为使分次 K-线性映射 \(f: \mathbf{M} \to \mathbf{E}\) （次数为 \(\delta'\) ）可以扩张为一个 \(\varepsilon'\) -导子
\end{enumerate}

\[
d: \bigwedge (\mathbf {M}) \rightarrow \mathbf {E}
\]

（次数为 \(\delta'\) ），必要且充分条件是：对每个齐次元素 \(x\) （属于 \(\mathbf{M}\) ），

\[
x. f (x) + \varepsilon_ {\delta ^ {\prime}, (\deg (x), 1)} ^ {\prime} f (x). x = 0.\tag{34}
\]

该 \(\varepsilon'\) -导子 \(d\) 则是唯一的。

第 2 号的注记 2 应用于如下情形：E 上的外部 B-模律之一（其中 B 等于 T(M)、S(M) 或 ∧(M)）被修改；如此修改后的外部律依 (1) （第 1 号）仍是一个 B-模律，且由此在 E 上得到的 B-模律与另一个 B-模结构仍然相容。这时只需将命题 14 应用于 A = K 与 \(d_{0} = 0\) 的情形。

例 (1). 在命题 14 的应用中，注意若 \(d_0 = 0\) ，条件 (29) 仅表示 \(d_1\) 是 A-线性的。若特别取 \(\mathbf{E} = \mathbf{B}\) 以及由 B 上的环结构导出的 (B, B)-双模结构，则当 \(d_1\) 取为某个自同态 \(s\) （ \(\mathbf{M}\) 的）与典范单射 \(\mathbf{M} \to \mathbf{B}\) 的复合时，条件 (30) 和 (31) 自动满足；这对 (30) 是显然的，因为 \(\mathbb{S}(\mathbf{M})\) 是交换的，而对 (31)，这由以下事实得出： \(x\) 和 \(s(x)\) 在 \(\bigwedge (\mathbf{M})\) 中都是次数 1 的。由此可以看出，每个自同态 \(s\) （ \(\mathbf{M}\) 的）都可以唯一地延拓为一个导子 \(\mathbf{D}_s\) （ \(\mathbb{T}(\mathbf{M})\) 的，相应地 \(\mathbb{S}(\mathbf{M})\) 的，相应地 \(\bigwedge (\mathbf{M})\) 的），其次数为 0。此外，对于两个自同态 \(s, t\) （ \(\mathbf{M}\) 的），

\[
[ \mathbf {D} _ {s}, \mathbf {D} _ {t} ] = \mathbf {D} _ {[ s, t ]}\tag{35}
\]

因为两边都是 \(\mathsf{T}(\mathbf{M})\) （相应地 \(\mathbb{S}(\mathbf{M})\) ，相应地 \(\bigwedge (\mathbf{M})\) ）的导子，且都等于 \([s, t]\) （在 \(\mathbf{M}\) 上）。

\(D_{s}\) 的表达式由第 5 号的公式 (21) 得到，该公式对 M 中的 \(x_{1}, x_{2}, \ldots, x_{n}\) 分别给出：

\[
\left\{ \begin{array}{l} \mathrm{D} _ {s} (x _ {1} \otimes x _ {2} \otimes \dots \otimes x _ {n}) \\ \qquad = \sum_ {i = 1} ^ {n} x _ {1} \otimes \dots \otimes x _ {i - 1} \otimes s (x _ {i}) \otimes x _ {i + 1} \otimes \dots \otimes x _ {n} \\ \mathrm{D} _ {s} (x _ {1} x _ {2} \dots x _ {n}) = \sum_ {i = 1} ^ {n} x _ {1} \dots x _ {i - 1} s (x _ {i}) x _ {i + 1} \dots x _ {n} \\ \mathrm{D} _ {s} (x _ {1} \wedge x _ {2} \wedge \dots \wedge x _ {n}) \\ \qquad = \sum_ {i = 1} ^ {n} x _ {1} \wedge \dots \wedge x _ {i - 1} \wedge s (x _ {i}) \wedge x _ {i + 1} \wedge \dots \wedge x _ {n}. \end{array} \right.\tag{36}
\]

在代数 \(\bigwedge (\mathbf{M})\) 的情形，对 \(\mathbf{D}_s\) 有如下的诠释：

命题 15. 若 M 是有限秩 n 的自由 K-模，则对 M 的每个自同态 s ，限制在 \(\bigwedge^{n}(\mathbf{M})\) 上时，导子 \(\mathbf{D}_{s}\) 是比率为 \(\operatorname{Tr}(s)\) 的位似。

设 \((e_j)_{1\leqslant j\leqslant n}\) 是 M 的一个基，并记 \(e = e_1\wedge e_2\wedge \dots \wedge e_n\) 。若

\[
s (e _ {j}) = \sum_ {k = 1} ^ {n} \alpha_ {j k} e _ {k},
\]

(36) 中的第三个公式给出

\[
\mathrm{D} _ {s} (e) = \sum_ {i = 1} ^ {n} e _ {1} \wedge \dots \wedge e _ {i - 1} \wedge s \left(e _ {i}\right) \wedge e _ {i + 1} \wedge \dots \wedge e _ {n} = \left(\sum_ {j = 1} ^ {n} \alpha_ {j j}\right) e.
\]

例 (2). 在命题 14 的推论的第 (iii) 部分中，令 \(\Delta = \{0\}\) ，此时 \(\bigwedge (\mathbf{M})\) 上的分次是通常的类型为 \(\mathbf{Z}\) 的分次；另一方面取 \(\varepsilon(p, q) = (-1)^{pq}\) 。那么，对每个线性形式 \(x^{*} \in \mathbf{M}^{*}\) （ \(\mathbf{M}\) 上的）， \(x \mapsto \langle x, x^{*} \rangle\) 是一个次数为 \(-1\) 的分次 K-线性映射（从 \(\mathbf{M}\) 到 \(\bigwedge (\mathbf{M})\) ），满足关系式 (34)；于是存在一个反导子 \(i(x^{*})\) （ \(\bigwedge (\mathbf{M})\) 上的，次数为 \(-1\) ），使得（依第 5 号的公式 (21)）

\[
i \left(x ^ {*}\right) \left(x _ {1} \wedge \dots \wedge x _ {n}\right)
\]

\[
= \sum_ {i = 1} ^ {n} (- 1) ^ {i - 1} \left\langle x _ {i}, x ^ {*} \right\rangle x _ {1} \wedge \dots \wedge x _ {i - 1} \wedge x _ {i + 1} \wedge \dots \wedge x _ {n}
\]

并且这是将在 § 11 第 9 号公式 (68) 中定义的内积的一个特例。

命题 16. 设 A 为交换 K-代数， \(\mathbf{M}_i\) ( \(1 \leqslant i \leqslant n\) ) 和 P 为 A-模，H 为从 \(\mathbf{M}_1 \times \mathbf{M}_2 \times \cdots \times \mathbf{M}_n\) 到 P 中的 A-多重线性映射所构成的 A-模。假设给定代数 A 的一个 K-导子 \(d_{0}:A\to A\) ，对每个 i 给定一个 K-线性映射 \(d_{i}:M_{i}\to M_{i}\) 以及一个 K-线性映射 D:P→P，使得对 \(1\leqslant i\leqslant n\) ， \((d_{0},d_{i},d_{i})\) 是 \((\mathbf{A},\mathbf{M}_{i},\mathbf{M}_{i})\) 到自身的一个 K-导子，且 \((d_{0},\mathbf{D},\mathbf{D})\) 是 \((\mathbf{A},\mathbf{P},\mathbf{P})\) 到自身的一个 K-导子。则存在一个 K-线性映射 \(D^{\prime}:H\to H\) ，使得 \((d_{0},D^{\prime},D^{\prime})\) 是 \((\mathbf{A},\mathbf{H},\mathbf{H})\) 到自身的一个 K-导子，且

\begin{enumerate}
\def\labelenumi{(\arabic{enumi})}
\setcounter{enumi}{36}
\tightlist
\item
  \(\mathbf{D}(f(x_1, \ldots, x_n))\)
\end{enumerate}

对所有 \(x_{i} \in \mathbf{M}_{i}\) （ \(1 \leqslant i \leqslant n\) ）和 \(f \in \mathbf{H}\) 成立。

我们证明，对 \(f \in H\) ，映射 \(D'f\) （从 \(M_{1} \times M_{2} \times \cdots \times M_{n}\) 到 P 中的，由 (37) 定义）是 A-多重线性的。对 \(a \in A\) ，

\[
\begin{array}{r l} (\mathrm{D} ^ {\prime} f) (a x _ {1}, x _ {2}, \ldots , x _ {n}) & = \mathrm{D} (a f (x _ {1}, \ldots , x _ {n})) - f (d _ {1} (a x _ {1}), x _ {2}, \ldots , x _ {n}) \\ & \quad - a \sum_ {i = 2} ^ {n} f (x _ {1}, \ldots , x _ {i - 1}, d _ {i} x _ {i}, x _ {i + 1}, \ldots , x _ {n}) \end{array}
\]

且根据假设

\[
\mathrm{D} (a f (x _ {1}, \dots , x _ {n})) = (d _ {0} a) f (x _ {1}, \dots , x _ {n}) + a \mathrm{D} (f (x _ {1}, \dots , x _ {n}))
\]

和 \(d_{1}(ax_{1}) = (d_{0}a)x_{1} + a.d_{1}x_{1}\) ，这给出

\[
(\mathrm{D} ^ {\prime} f) (a x _ {1}, x _ {2}, \dots , x _ {n}) = a. (\mathrm{D} ^ {\prime} f) (x _ {1}, \dots , x _ {n})
\]

而关于每个 \(x_{i}\) 的线性性可类似地证明。另一方面，

\[
\begin{array}{r l} (\mathrm{D} ^ {\prime} (a f)) (x _ {1}, \dots , x _ {n}) & = \mathrm{D} (a f (x _ {1}, \dots , x _ {n})) \\ & - \sum_ {i = 1} ^ {n} a f (x _ {1}, \dots , x _ {i - 1}, d _ {i} x _ {i}, x _ {i + 1}, \dots , x _ {n}) \\ & = (d _ {0} a) f (x _ {1}, \dots , x _ {n})) + a \mathrm{D} (f (x _ {1}, \dots , x _ {n})) \\ & - \sum_ {i = 1} ^ {n} a f (x _ {1}, \dots , x _ {i - 1}, d _ {i} x _ {i}, x _ {i + 1}, \dots , x _ {n}) \\ & = (d _ {0} a) f (x _ {1}, \dots , x _ {m}) + a (\mathrm{D} ^ {\prime} f) (x _ {1}, \dots , x _ {m}) \end{array}
\]

换句话说

\[
\mathbf {D} ^ {\prime} (a f) = \left(d _ {0} a\right) f + a (\mathbf {D} ^ {\prime} f)
\]

这便确立了该命题。

例. (3) 将命题 16 应用于情形 \(n = 1\) ， \(\mathbf{M}_1 = \mathbf{M}\) ， \(P = A\) ，则 \(H = M^*\) ，即 \(M\) 的对偶，并且可以看出，由 K-导子 \((d_0, d, d)\) （ \((A, M, M)\) 的）可导出一个 K-导子 \((d_0, d^*, d^*)\) （ \((A, M^*, M^*)\) 的），使得

\[
d _ {0} \langle m, m ^ {*} \rangle = \langle d m, m ^ {*} \rangle + \langle m, d ^ {*} m ^ {*} \rangle\tag{38}
\]

对 \(m \in \mathbf{M}\) 和 \(m^* \in \mathbf{M}^*\) 成立。从 \(\mathbf{M} \oplus \mathbf{M}^*\) 到自身的、等于 \(d\) （在 \(\mathbf{M}\) 上）且等于 \(d^{*}\) （在 \(\mathbf{M}^{*}\) 上）的那个 K-线性映射于是满足条件 (29)，因此存在一个 K-导子 \(\mathbf{D}\) （A-代数 \(\mathsf{T}(\mathbf{M} \oplus \mathbf{M}^{*})\) 的），它约化为 \(d_{0}\) （在 \(\mathbf{A}\) 上）、为 \(d\) （在 \(\mathbf{M}\) 上）、为 \(d^{*}\) （在 \(\mathbf{M}^{*}\) 上）。限制 \(d_{\mathbf{J}}^{\mathbf{I}}\) ——即 \(\mathbf{D}\) 在子 A-模 \(\mathsf{T}_{\mathbf{J}}^{\mathbf{I}}(\mathbf{M})\) （ \(\mathsf{T}(\mathbf{M} \oplus \mathbf{M}^{*})\) 中的； \(\S 5\) ，第 6 号）上的限制——是 \(\mathsf{T}_{\mathbf{J}}^{\mathbf{I}}(\mathbf{M})\) 的一个 K-自同态，使得 \((d_{0}, d_{\mathbf{J}}^{\mathbf{I}}, d_{\mathbf{J}}^{\mathbf{I}})\) 是 \((\mathbf{A}, \mathsf{T}_{\mathbf{J}}^{\mathbf{I}}(\mathbf{M}), \mathsf{T}_{\mathbf{J}}^{\mathbf{I}}(\mathbf{M}))\) 的一个 K-导子。此外，对于 \(i \in \mathbf{I}\) ， \(j \in \mathbf{J}\) ，若记 \(\mathbf{I}' = \mathbf{I} - \{i\}\) ， \(\mathbf{J}' = \mathbf{J} - \{j\}\) ，则立即可验证，对收缩 \(c_{j}^{t}\) （ \(\S 5\) ，第 6 号）

\[
c _ {j} ^ {t} (d _ {J} ^ {\mathrm{I}} (z)) = d _ {J ^ {\prime}} ^ {\mathrm{I} ^ {\prime}} (c _ {j} ^ {t} (z)) \quad \text { for   all } z \in \mathbf {T} _ {J} ^ {\mathrm{I}} (\mathbf {M}).
\]

\begin{enumerate}
\def\labelenumi{(\arabic{enumi})}
\setcounter{enumi}{3}
\tightlist
\item
  设 \(M_{i}\) ( \(1 \leqslant i \leqslant 3\) ) 是三个 A-模，并且对每个 i 假设 \((d_{0}, d_{i}, d_{i})\) 是 \((\mathbf{A}, \mathbf{M}_{i}, \mathbf{M}_{i})\) 的一个导子；再对 n = 1 应用命题 16，则可导出一个导子 \((d_{0}, d_{ij}, d_{ij})\) （ \((\mathbf{A}, \mathbf{H}_{ij}, \mathbf{H}_{ij})\) 的，其中 \(\mathrm{H}_{ij} = \mathrm{Hom}_{\mathbf{A}}(\mathbf{M}_{i}, \mathbf{M}_{j})\) ），对每个有序对 \((i, j)\) 都是如此。利用这个记号，对 \(u \in \mathrm{Hom}_{\mathbf{A}}(\mathbf{M}_{1}, \mathbf{M}_{2})\) 和 \(v \in \mathrm{Hom}_{\mathbf{A}}(\mathbf{M}_{2}, \mathbf{M}_{3})\) ，
\end{enumerate}

\[
d _ {1 3} (v \circ u) = (d _ {2 3} v) \circ u + v \circ (d _ {1 2} u)\tag{39}
\]

正如从定义出发立即验证的那样。

